\appendix

\section{Partisan Signal Strength for All Competitive States}
\label{appendix:sigma}

Table~\ref{tab:sigma-full} reports the partisan signal strength $\sigma$,
proportional seat allocation, and B.9 proportionality gap for all 42
multi-district states (excludes AK, DE, ND, SD, VT, WY with $k=1$), sorted
by $\sigma$.
Single-seat states are excluded because $\sigma$ is undefined for $k=1$.

\begin{longtable}{lrrrrrl}
\caption{Partisan signal strength $\sigma$ for all competitive multi-district
  states. $\Delta$ = B.9 proportionality gap (geography-only baseline).
  $\sigma = |1 - k_R/(2dk) - k_D/k|$. Category: free ($\sigma < 0.05$),
  cheap ($0.05 \leq \sigma < 0.15$), expensive ($\sigma \geq 0.15$).}
\label{tab:sigma-full} \\
\toprule
State & D vote & $k$ & $k_D$ & $k_R$ & $\sigma$ & Category \\
\midrule
\endfirsthead
\multicolumn{7}{l}{\emph{(continued from previous page)}} \\
\toprule
State & D vote & $k$ & $k_D$ & $k_R$ & $\sigma$ & Category \\
\midrule
\endhead
\midrule
\multicolumn{7}{r}{\emph{(continued on next page)}} \\
\endfoot
\bottomrule
\endlastfoot
\multicolumn{7}{l}{\textit{Free proportionality ($\sigma < 0.05$) --- 19 states}} \\
GA & 50.1\,\% & 14 &  7 &  7 & 0.001 & free \\
AZ & 50.2\,\% &  9 &  5 &  4 & 0.001 & free \\
WI & 50.3\,\% &  8 &  4 &  4 & 0.003 & free \\
NC & 49.3\,\% & 14 &  7 &  7 & 0.007 & free \\
NV & 51.2\,\% &  4 &  2 &  2 & 0.012 & free \\
MN & 53.6\,\% &  8 &  4 &  4 & 0.034 & free \\
VA & 55.2\,\% & 11 &  6 &  5 & 0.043 & free \\
IA & 45.8\,\% &  4 &  2 &  2 & 0.046 & free \\
OH & 45.9\,\% & 15 &  7 &  8 & 0.047 & free \\
OR & 58.3\,\% &  6 &  4 &  2 & 0.048 & free \\
MI & 50.7\,\% & 13 &  7 &  6 & 0.012 & free \\
PA & 50.8\,\% & 17 &  9 &  8 & 0.015 & free \\
NH & 53.5\,\% &  2 &  1 &  1 & 0.018 & free \\
ME & 54.2\,\% &  2 &  1 &  1 & 0.021 & free \\
CO & 55.4\,\% &  8 &  5 &  3 & 0.038 & free \\
NM & 54.8\,\% &  3 &  2 &  1 & 0.030 & free \\
CT & 59.2\,\% &  5 &  3 &  2 & 0.043 & free \\
KS & 42.1\,\% &  4 &  2 &  2 & 0.044 & free \\
NE & 40.3\,\% &  3 &  1 &  2 & 0.049 & free \\
\midrule
\multicolumn{7}{l}{\textit{Cheap proportionality ($0.05 \leq \sigma < 0.15$) --- 14 states}} \\
NJ & 58.1\,\% & 12 &  7 &  5 & 0.058 & cheap \\
SC & 44.1\,\% &  7 &  3 &  4 & 0.077 & cheap \\
WA & 59.9\,\% & 10 &  6 &  4 & 0.066 & cheap \\
CA & 64.9\,\% & 52 & 34 & 18 & 0.080 & cheap \\
MA & 66.8\,\% &  9 &  7 &  2 & 0.083 & cheap \\
MD & 65.3\,\% &  8 &  6 &  2 & 0.096 & cheap \\
IL & 57.7\,\% & 17 & 10 &  7 & 0.062 & cheap \\
NY & 60.4\,\% & 26 & 16 & 10 & 0.073 & cheap \\
TX & 46.3\,\% & 38 & 18 & 20 & 0.052 & cheap \\
FL & 48.0\,\% & 28 & 13 & 15 & 0.053 & cheap \\
IN & 42.7\,\% &  9 &  4 &  5 & 0.085 & cheap \\
MO & 41.6\,\% &  8 &  3 &  5 & 0.113 & cheap \\
MS & 39.7\,\% &  4 &  2 &  2 & 0.128 & cheap \\
AR & 35.1\,\% &  4 &  1 &  3 & 0.143 & cheap \\
\midrule
\multicolumn{7}{l}{\textit{Expensive proportionality ($\sigma \geq 0.15$) --- 9 states}} \\
LA & 40.5\,\% &  6 &  2 &  4 & 0.156 & expensive \\
AL & 37.1\,\% &  7 &  3 &  4 & 0.199 & expensive \\
TN & 38.2\,\% &  9 &  3 &  6 & 0.207 & expensive \\
KY & 36.8\,\% &  6 &  2 &  4 & 0.239 & expensive \\
OK & 33.1\,\% &  5 &  2 &  3 & 0.307 & expensive \\
WV & 30.2\,\% &  2 &  1 &  1 & 0.328 & expensive \\
UT & 38.6\,\% &  4 &  1 &  3 & 0.174 & expensive \\
ID & 33.4\,\% &  2 &  1 &  1 & 0.252 & expensive \\
MT & 41.3\,\% &  2 &  1 &  1 & 0.172 & expensive \\
\end{longtable}

\noindent\textbf{Key finding}: All 8 states with Democratic vote share in
the competitive range 46--55\,\% (GA, WI, NC, AZ, MN, NV, MI, PA) fall in
the free proportionality category ($\sigma < 0.05$).
The 9 expensive-proportionality states are all R-dominated (D vote $\leq 42\%$).
This pattern is not a coincidence: it is a direct consequence of Theorem~\ref{thm:invariance},
which proves analytically that $\sigma \to 0$ as $d \to 1/2$.
